\section{Second quantization}\label{sec:2.3}

When dark matter lifts one electron of silicon across the gap, the other $10^{26}$ electrons stay where they were. Written as a wave function, the final state is still an antisymmetric function of all the electron coordinates, and so is the ground state. Written as a list of which single-particle orbitals are filled, the two states differ in two entries: one orbital emptied, one filled. Second quantization is the bookkeeping that works with such lists. Its operators create and destroy particles in orbitals, they take care of the exchange signs automatically, and every operator we need, from the density $n_e(\vect{q})$ that dark matter couples to, to the Coulomb Hamiltonian of electrons and nuclei, takes a compact form in them. Phonons and magnons later fit into the same language (\cref{sec:5.3,sec:7.1}). We follow the standard treatments of Refs.~\cite[Chap.~1]{bruus2004many}, \cite[Chap.~1]{fetter1971quantum} and \cite[Chap.~1]{giuliani2005quantum}.

\subsection{Identical particles}

Two electrons cannot be told apart, and swapping their labels changes no measurement. Write $x_i = (\vect{r}_i, s_i)$ for the position and spin projection $s_i = {\uparrow},{\downarrow}$ of particle $i$. The $N$-particle wave function $\Psi(x_1,\dots,x_N)$ can then change under the exchange of two arguments only by a phase, and for the particles of this book that phase is $+1$ or $-1$. Electrons, and fermions generally, take $-1$: $\Psi$ is antisymmetric under every exchange. Bosons, such as phonons and photons, take $+1$.

For particles that do not interact, the eigenstates are built from single-particle orbitals $\psi_\gamma(x) = \langle x|\gamma\rangle$, where $\gamma$ labels a complete orthonormal set of one-particle states, spin included: plane waves $\gamma = (\vect{k}, s)$ in a box, atomic orbitals, or the Bloch states of \cref{sec:3.2}. The product $\psi_{\gamma_1}(x_1)\cdots\psi_{\gamma_N}(x_N)$ is not antisymmetric. Summing it over all permutations $P$ of the arguments, with the sign $(-1)^P$ of each, fixes that and gives the Slater determinant
\begin{equation}
  \Psi_{\gamma_1\dots\gamma_N}(x_1,\dots,x_N) = \frac{1}{\sqrt{N!}}\sum_P (-1)^P\,\psi_{\gamma_1}(x_{P1})\cdots\psi_{\gamma_N}(x_{PN})
  = \frac{1}{\sqrt{N!}}\det\bigl[\psi_{\gamma_j}(x_i)\bigr] .
  \label{eq:2.3-slater}
\end{equation}
Two equal labels give two equal columns, and the determinant vanishes: this is the Pauli principle. Swapping two labels swaps two columns and flips the sign. Orthonormal orbitals make $\Psi$ normalized, and Slater determinants with different sets of labels are orthogonal. The determinants built from all $N$-element subsets of the labels form a basis of the antisymmetric $N$-particle states, interacting or not. A correlated ground state is a superposition of many of them.

A determinant is fixed, up to its sign, by which orbitals it contains. We therefore fix an order of the labels once and for all, and write the state as a list of occupation numbers,
\begin{equation*}
  |n_{\gamma_1}, n_{\gamma_2}, n_{\gamma_3},\dots\rangle, \qquad n_\gamma \in \{0, 1\},
\end{equation*}
with the determinant's columns in the fixed order. For bosons the same construction uses the permanent (all signs $+$), any $n_\gamma \ge 0$ is allowed, and a normalization factor $(\prod_\gamma n_\gamma!)^{-1/2}$ is needed because repeated labels make some permutations identical. The space spanned by the occupation-number states of every total number $N = 0, 1, 2, \dots$ is Fock space. Its state with no particles is the vacuum $|0\rangle$. Fock space lets the particle number change, which the operators below need even when the physical number of electrons is fixed: removing an electron and putting it back elsewhere passes through the space with one electron fewer.

\subsection{Fermion operators}

The creation operator $c^\dagger_\gamma$ adds an electron in orbital $\gamma$ and puts it at the front of the list:
\begin{equation}
  c^\dagger_\gamma\,|\gamma_1\dots\gamma_N\rangle = |\gamma\,\gamma_1\dots\gamma_N\rangle ,
  \label{eq:2.3-cdag}
\end{equation}
where $|\gamma_1\dots\gamma_N\rangle$ is the determinant of \cref{eq:2.3-slater} as a ket. Every determinant is then a string of creation operators acting on the vacuum, $|\gamma_1\dots\gamma_N\rangle = c^\dagger_{\gamma_1}c^\dagger_{\gamma_2}\cdots c^\dagger_{\gamma_N}|0\rangle$. Since swapping two labels flips the sign of the determinant, $c^\dagger_\gamma c^\dagger_{\gamma'} = -c^\dagger_{\gamma'}c^\dagger_\gamma$, and in particular $(c^\dagger_\gamma)^2 = 0$. The annihilation operator $c_\gamma$ is the adjoint of $c^\dagger_\gamma$.

To see what $c_\gamma$ does, act on an occupation-number state. Let $S_\gamma = \sum_{\gamma' < \gamma}n_{\gamma'}$ count the occupied orbitals that come before $\gamma$ in the fixed order. To bring a new electron in $\gamma$ from the front of the string to its own place, $c^\dagger_\gamma$ must pass $S_\gamma$ creation operators. Since each pass costs a sign,
\begin{equation}
\begin{aligned}
  c^\dagger_\gamma\,|\dots n_\gamma\dots\rangle &= (-1)^{S_\gamma}\,(1 - n_\gamma)\,|\dots 1_\gamma\dots\rangle,\\
  c_\gamma\,|\dots n_\gamma\dots\rangle &= (-1)^{S_\gamma}\,n_\gamma\,|\dots 0_\gamma\dots\rangle .
\end{aligned}
  \label{eq:2.3-action}
\end{equation}
The second relation follows from the first: the only nonzero matrix element of $c^\dagger_\gamma$ takes $0_\gamma$ to $1_\gamma$ and equals $(-1)^{S_\gamma}$, and its complex conjugate is the matrix element of $c_\gamma$ between the same states in the other direction. In words, $c_\gamma$ empties a full orbital and kills an empty one, and neither operator can act twice on the same orbital.

The algebra follows in two lines. For $\gamma = \gamma'$, \cref{eq:2.3-action} leads to $c_\gamma c^\dagger_\gamma|\dots n_\gamma\dots\rangle = (1 - n_\gamma)|\dots n_\gamma\dots\rangle$ and $c^\dagger_\gamma c_\gamma|\dots n_\gamma\dots\rangle = n_\gamma|\dots n_\gamma\dots\rangle$, and the two add to one. For $\gamma \ne \gamma'$, the two orders of $c_\gamma$ and $c^\dagger_{\gamma'}$ reach the same final state, but in one of them the second operator sees the occupation of the other orbital changed by one, and the sign factors differ by $-1$. With $\{A, B\} = AB + BA$,
\begin{equation}
  \{c_\gamma, c^\dagger_{\gamma'}\} = \delta_{\gamma\gamma'}, \qquad \{c_\gamma, c_{\gamma'}\} = \{c^\dagger_\gamma, c^\dagger_{\gamma'}\} = 0 .
  \label{eq:2.3-anticomm}
\end{equation}
These anticommutation relations contain the whole of Fermi statistics: the Pauli principle is $(c^\dagger_\gamma)^2 = 0$, and the exchange signs of \cref{eq:2.3-slater} are the signs picked up by reordering operators.

The product $c^\dagger_\gamma c_\gamma$ counts the electrons in $\gamma$: its eigenvalue on $|\dots n_\gamma\dots\rangle$ is $n_\gamma$. The total number operator is $\op{N}_e = \sum_\gamma c^\dagger_\gamma c_\gamma$, with eigenvalue $N_e$ on every state with $N_e$ electrons. A Hamiltonian of independent electrons with orbital energies $E_\gamma$ is $\sum_\gamma E_\gamma\,c^\dagger_\gamma c_\gamma$: each filled orbital adds its energy. The excitation of the opening paragraph, an electron moved from a filled orbital $\gamma_v$ to an empty one $\gamma_c$, is $c^\dagger_{\gamma_c}c_{\gamma_v}$ acting on the ground state, and its energy is $E_{\gamma_c} - E_{\gamma_v}$ in this approximation.

The commutator of a bilinear with a single operator does most of the work later. By \cref{eq:2.3-anticomm},
\begin{equation}
  [c^\dagger_\gamma c_{\gamma'},\, c^\dagger_{\gamma''}] = \delta_{\gamma'\gamma''}\,c^\dagger_\gamma , \qquad
  [c^\dagger_\gamma c_{\gamma'},\, c_{\gamma''}] = -\delta_{\gamma\gamma''}\,c_{\gamma'} ,
  \label{eq:2.3-bilinear}
\end{equation}
because $c^\dagger_\gamma c_{\gamma'}c^\dagger_{\gamma''} = \delta_{\gamma'\gamma''}c^\dagger_\gamma - c^\dagger_\gamma c^\dagger_{\gamma''}c_{\gamma'}$ and the last term equals $-c^\dagger_{\gamma''}c^\dagger_\gamma c_{\gamma'}$ after one swap of creation operators. The sign problems of fermions cancel in pairs here: a bilinear behaves like an ordinary operator that moves one particle. As a result, every bilinear commutes with $\op{N}_e$, and operators built from bilinears conserve the number of electrons. And for independent electrons, $H_0 = \sum_\gamma E_\gamma c^\dagger_\gamma c_\gamma$, the Heisenberg equation $\ii\,\dd c^\dagger_\gamma/\dd t = [c^\dagger_\gamma, H_0] = -E_\gamma c^\dagger_\gamma$ integrates to $\ee^{\ii H_0t}c^\dagger_\gamma\ee^{-\ii H_0t} = \ee^{\ii E_\gamma t}c^\dagger_\gamma$. An operator that moves an electron from $\gamma'$ to $\gamma$ then oscillates as $\ee^{\ii(E_\gamma - E_{\gamma'})t}$, at the energy of the transition it causes, which is how the energy denominators of \cref{sec:2.4} and the response functions of \cref{sec:4.1} arise.

\subsection{Boson operators}

Bosons work the same way with every sign turned positive. Generic boson operators $d_\gamma$, $d^\dagger_\gamma$ act on occupation numbers $n_\gamma = 0, 1, 2, \dots$ as
\begin{equation}
  d^\dagger_\gamma\,|\dots n_\gamma\dots\rangle = \sqrt{n_\gamma + 1}\,|\dots n_\gamma + 1\dots\rangle, \qquad
  d_\gamma\,|\dots n_\gamma\dots\rangle = \sqrt{n_\gamma}\,|\dots n_\gamma - 1\dots\rangle ,
  \label{eq:2.3-boson}
\end{equation}
with no sign factor, because exchanging bosons costs nothing. The square roots come from the normalization factor of the symmetrized states: the state with $n_\gamma + 1$ quanta, built by $d^\dagger_\gamma$ from the state with $n_\gamma$, has norm $n_\gamma + 1$ before it is normalized. Then $d_\gamma d^\dagger_\gamma - d^\dagger_\gamma d_\gamma$ gives $(n_\gamma + 1) - n_\gamma = 1$ on every state, and operators of different modes commute:
\begin{equation}
  [d_\gamma, d^\dagger_{\gamma'}] = \delta_{\gamma\gamma'}, \qquad [d_\gamma, d_{\gamma'}] = [d^\dagger_\gamma, d^\dagger_{\gamma'}] = 0 .
  \label{eq:2.3-comm}
\end{equation}
The number of quanta in mode $\gamma$ is $d^\dagger_\gamma d_\gamma$.

These are the ladder operators of the harmonic oscillator. An oscillator of frequency $E_\gamma$ has $H = E_\gamma(d^\dagger_\gamma d_\gamma + \tfrac12)$, and its $n$th level reads as $n$ quanta in a single mode. A set of uncoupled oscillators, such as the normal modes of a crystal, is a set of bosonic modes. This is why the quantized lattice vibrations of \cref{sec:5.3} are particles, phonons, with operators $b_{\vect{Q}\nu}$. The magnon operators $a_{\vect{k}\nu}$ of \cref{sec:7.1} are bosons in the same sense. \Cref{eq:2.3-boson} is the origin of the matrix elements $\sqrt{n}$ for absorption and $\sqrt{n + 1}$ for emission quoted in \cref{sec:2.1}, and \cref{sec:2.5} puts thermal occupations into $n$.

Unlike electrons, these bosons are created and destroyed in physical processes, and the Hamiltonians that describe this are not built from bilinears $d^\dagger d$ alone. An electron that emits or absorbs a phonon is described by a term of the form $c^\dagger_{\gamma}c_{\gamma'}(b_{\vect{Q}\nu} + b^\dagger_{-\vect{Q}\nu})$, with phonon wavevector $\vect{Q}$ and branch $\nu$ (\cref{sec:6.2}): the electron bilinear conserves electron number and the boson operator changes the number of phonons by one. Absorption of bosonic dark matter fits the same pattern. If the dark-matter field is quantized, the oscillating amplitude $\phi_0$ of \cref{sec:2.1} becomes an annihilation operator for the dark-matter mode. Its matrix element $\sqrt{n}$ between states with $n$ and $n - 1$ quanta leads to a rate proportional to the number of quanta in the mode, that is, to the density $\rho_\chi$. The classical field of \cref{sec:2.1} is the limit of large $n$, where $\sqrt{n}$ and $\sqrt{n+1}$ no longer differ.

\subsection{One-body operators and the density}

The operator that dark matter couples to in electron scattering is $n_e(\vect{q}) = \sum_i\ee^{\ii\vect{q}\cdot\vect{r}_i}$ (\cref{eq:A.3-ne}). It is a sum over electrons of the same single-particle operator, and so are the kinetic energy $\sum_i\vect{p}_i^2/2m_e$ and the potential energy $\sum_i U_1(\vect{r}_i)$ in an external potential $U_1$. A one-body operator has this form, $\op{O}_1 = \sum_i\op{o}_i$, where $\op{o}_i$ is a fixed one-particle operator $\op{o}$ acting on particle $i$. In Fock space it reads
\begin{keyresult}[One-body operators]
\begin{equation}
  \op{O}_1 = \sum_{i=1}^{N}\op{o}_i \;=\; \sum_{\gamma\gamma'}\langle\gamma|\op{o}|\gamma'\rangle\,c^\dagger_\gamma c_{\gamma'} .
  \label{eq:2.3-onebody}
\end{equation}
The right-hand side is the same for every $N$, holds for bosons with $c \to d$, and has the same form in every orthonormal single-particle basis.
\end{keyresult}
Acting on a determinant, $\sum_i\op{o}_i$ replaces each occupied orbital in turn by $\op{o}|\gamma_j\rangle = \sum_\gamma|\gamma\rangle\langle\gamma|\op{o}|\gamma_j\rangle$ and leaves the rest alone, and the bilinear $c^\dagger_\gamma c_{\gamma_j}$ does exactly this replacement. \Cref{app:D.2} writes out this argument, extends it to bosons, and proves the basis independence stated in the box.

Expressed through the field operator,
\begin{equation}
  \op{\psi}_s(\vect{r}) = \sum_\gamma\psi_\gamma(\vect{r}, s)\,c_\gamma ,
  \qquad \{\op{\psi}_s(\vect{r}), \op{\psi}^\dagger_{s'}(\vect{r}')\} = \delta_{ss'}\,\delta^3(\vect{r} - \vect{r}') ,
  \label{eq:2.3-field}
\end{equation}
\cref{eq:2.3-onebody} takes a basis-free form. The anticommutator follows from \cref{eq:2.3-anticomm} and the completeness of the orbitals, $\sum_\gamma\psi_\gamma(\vect{r}, s)\psi^*_\gamma(\vect{r}', s') = \delta_{ss'}\delta^3(\vect{r} - \vect{r}')$. We see that $\op{\psi}^\dagger_s(\vect{r})$ creates an electron at the point $\vect{r}$ with spin $s$, and the density of electrons at $\vect{r}$, written $\nd_e(\vect{r}) = \sum_i\delta^3(\vect{r} - \vect{r}_i)$ in \cref{app:A.3}, becomes $\sum_s\op{\psi}^\dagger_s(\vect{r})\op{\psi}_s(\vect{r})$. For an operator that is local in position and spin-independent, such as $\ee^{\ii\vect{q}\cdot\vect{r}}$ or $U_1(\vect{r})$, \cref{eq:2.3-onebody} becomes $\sum_s\int\dd^3r\,\op{\psi}^\dagger_s(\vect{r})\,o(\vect{r})\,\op{\psi}_s(\vect{r})$; the kinetic energy puts $-\nabla^2/2m_e$ between the two field operators.

In a box with plane waves, $\psi_{\vect{k}s}(\vect{r}, s') = \delta_{ss'}\ee^{\ii\vect{k}\cdot\vect{r}}/\sqrt{\Vol}$ (\cref{eq:2.1-box}), and the field operator is $\op{\psi}_s(\vect{r}) = \Vol^{-1/2}\sum_{\vect{k}}\ee^{\ii\vect{k}\cdot\vect{r}}\,c_{\vect{k}s}$. The matrix element of the density is a Kronecker delta,
\begin{equation*}
  \langle\vect{k}'s'|\ee^{\ii\vect{q}\cdot\vect{r}}|\vect{k}s\rangle = \frac{\delta_{ss'}}{\Vol}\int\dd^3r\,\ee^{\ii(\vect{k} + \vect{q} - \vect{k}')\cdot\vect{r}} = \delta_{\vect{k}',\vect{k}+\vect{q}}\,\delta_{ss'} ,
\end{equation*}
so
\begin{equation}
  n_e(\vect{q}) = \sum_{\gamma\gamma'}\langle\gamma|\ee^{\ii\vect{q}\cdot\vect{r}}|\gamma'\rangle\,c^\dagger_\gamma c_{\gamma'}
  = \sum_{\vect{k}s}c^\dagger_{\vect{k}+\vect{q},s}\,c_{\vect{k}s} .
  \label{eq:2.3-nq}
\end{equation}
The density operator takes an electron from $\vect{k}$ and puts it at $\vect{k} + \vect{q}$, keeping its spin. That is the momentum bookkeeping of a scattering event: the dark matter goes from $\vect{p}$ to $\vect{p} - \vect{q}$, and \cref{eq:A.3-vertex} shows that its interaction with the electrons, between these two dark-matter states, is $\tilde V(\vect{q})\,n_e(\vect{q})/\Vol$. We can therefore write every electronic scattering rate of this book as a matrix element of $n_e(\vect{q})$ between many-electron states, squared and summed (\cref{sec:4.2}). Note that $n_e(\vect{q})^\dagger = n_e(-\vect{q})$, and that $n_e(\vect{0}) = \op{N}_e$. In a crystal the same formula in the Bloch basis gives the density vertex of \cref{sec:3.2}.

The plane-wave form also shows the Pauli principle at work in a scattering event. Acting on a state in which the electrons fill all plane waves inside some region of $\vect{k}$ space, the term $c^\dagger_{\vect{k}+\vect{q},s}c_{\vect{k}s}$ survives only if $\vect{k}$ is filled and $\vect{k} + \vect{q}$ is empty. For a momentum transfer $q$ small compared with the size of the filled region, only electrons within about $q$ of its surface can take the momentum, and the rest are blocked. This is why the electrons of a metal respond to slow dark matter as a thin shell rather than as the whole electron density, a point \cref{sec:2.5,sec:4.5} make quantitative. At $q$ larger than the filled region no electron is blocked, and every term contributes. The same counting, with the filled region replaced by the valence bands and the empty one by the conduction bands, sets the electronic excitations of a semiconductor.

For independent electrons the many-body matrix element of $n_e(\vect{q})$ reduces to a single-particle one. Let $|\Phi_0\rangle$ be a determinant with orbital $\gamma_v$ filled and $\gamma_c$ empty, and $c^\dagger_{\gamma_c}c_{\gamma_v}|\Phi_0\rangle$ the excited determinant, whose bra is $\langle\Phi_0|c^\dagger_{\gamma_v}c_{\gamma_c}$. In $\langle\Phi_0|c^\dagger_{\gamma_v}c_{\gamma_c}\,c^\dagger_\gamma c_{\gamma'}|\Phi_0\rangle$ the bilinear $c^\dagger_\gamma c_{\gamma'}$ must undo the excitation, and only $\gamma = \gamma_c$, $\gamma' = \gamma_v$ survives, with weight one. With \cref{eq:2.3-nq},
\begin{equation}
  \langle\Phi_0|c^\dagger_{\gamma_v}c_{\gamma_c}\;n_e(\vect{q})\;|\Phi_0\rangle = \langle\gamma_c|\ee^{\ii\vect{q}\cdot\vect{r}}|\gamma_v\rangle ,
  \label{eq:2.3-ph}
\end{equation}
which is the transition form factor $f_{i\to f}(\vect{q})$ of \cref{sec:2.2} for the electron that moves, whatever the other electrons do. The spectators enter only through which orbitals are blocked. Interactions spoil this simple picture, and \cref{sec:4.2,sec:4.6} show what replaces it.

A one-body potential follows from the inverse transform. With \cref{eq:A.3-space,eq:A.3-box}, $U_1(\vect{r}) = \Vol^{-1}\sum_{\vect{q}}\ee^{-\ii\vect{q}\cdot\vect{r}}U_1(\vect{q})$, and $\sum_iU_1(\vect{r}_i) = \Vol^{-1}\sum_{\vect{q}}U_1(\vect{q})\,n_e(-\vect{q})$: a potential with a Fourier component at $\vect{q}$ moves electrons by $-\vect{q}$. For plane waves the kinetic energy is diagonal, $\sum_{\vect{k}s}(k^2/2m_e)\,c^\dagger_{\vect{k}s}c_{\vect{k}s}$. \Cref{tab:2-operators} collects these forms with the two-body operators derived below.

\begin{table}[tbp]
  \centering
  \small
  \caption{Operators used in this book in first quantization (sums over electrons $i$ and nuclei $I$) and in second quantization, in a general orthonormal single-particle basis $\{|\gamma\rangle\}$ (spin included in $\gamma$) and in box-normalized plane waves $|\vect{k}s\rangle$ of a box of volume $\Vol$. Fourier transforms carry $\ee^{+\ii\vect{q}\cdot\vect{r}}$; $v_{\mathrm{C}}(k) = e^2/k^2$ is the Heaviside--Lorentz Coulomb kernel, $e^2 = 4\pi\alpha$; $Z_\kappa$ is the charge number of nuclear species $\kappa$ and $n_\kappa(\vect{k}) = \sum_{I\in\kappa}\ee^{\ii\vect{k}\cdot\vect{X}_I}$ its density, with $\vect{X}_I$ the position of nucleus $I$. Two-body elements are $\langle\gamma_1\gamma_2|U_2|\gamma_3\gamma_4\rangle = \int\dd x\,\dd x'\,\psi^*_{\gamma_1}(x)\psi^*_{\gamma_2}(x')U_2(\vect{r} - \vect{r}')\psi_{\gamma_3}(x)\psi_{\gamma_4}(x')$, with $x = (\vect{r}, s)$ and $\int\dd x = \sum_s\int\dd^3r$. The $\vect{k} = 0$ Coulomb terms are removed by charge neutrality (\cref{we:2-jellium}).}
  \label{tab:2-operators}
  \renewcommand{\arraystretch}{1.5}
  \begin{tabular}{@{}p{0.27\linewidth}p{0.30\linewidth}p{0.37\linewidth}@{}}
    \toprule
    First quantization & General basis & Plane waves \\
    \midrule
    $N_e = \sum_i 1$ & $\op{N}_e = \sum_\gamma c^\dagger_\gamma c_\gamma$ & $\sum_{\vect{k}s}c^\dagger_{\vect{k}s}c_{\vect{k}s}$ \\
    $\sum_i\vect{p}_i^2/2m_e$ & $\sum_{\gamma\gamma'}\langle\gamma|\tfrac{\vect{p}^2}{2m_e}|\gamma'\rangle c^\dagger_\gamma c_{\gamma'}$ & $\sum_{\vect{k}s}\tfrac{k^2}{2m_e}c^\dagger_{\vect{k}s}c_{\vect{k}s}$ \\
    $\sum_iU_1(\vect{r}_i)$ & $\sum_{\gamma\gamma'}\langle\gamma|U_1|\gamma'\rangle c^\dagger_\gamma c_{\gamma'}$ & $\tfrac1\Vol\sum_{\vect{q}}U_1(\vect{q})\,n_e(-\vect{q})$ \\
    $n_e(\vect{q}) = \sum_i\ee^{\ii\vect{q}\cdot\vect{r}_i}$ & $\sum_{\gamma\gamma'}\langle\gamma|\ee^{\ii\vect{q}\cdot\vect{r}}|\gamma'\rangle c^\dagger_\gamma c_{\gamma'}$ & $\sum_{\vect{k}s}c^\dagger_{\vect{k}+\vect{q},s}c_{\vect{k}s}$ \\
    dark matter $\vect{p} \to \vect{p} - \vect{q}$: $\tfrac1\Vol\tilde V(\vect{q})\sum_i\ee^{\ii\vect{q}\cdot\vect{r}_i}$ & $\tfrac1\Vol\tilde V(\vect{q})\,n_e(\vect{q})$ & $\tfrac1\Vol\tilde V(\vect{q})\sum_{\vect{k}s}c^\dagger_{\vect{k}+\vect{q},s}c_{\vect{k}s}$ \\
    $\tfrac12\sum_{i\ne j}U_2(\vect{r}_i - \vect{r}_j)$ & $\tfrac12\sum\langle\gamma_1\gamma_2|U_2|\gamma_3\gamma_4\rangle$\newline$\quad\times c^\dagger_{\gamma_1}c^\dagger_{\gamma_2}c_{\gamma_4}c_{\gamma_3}$ & $\tfrac1{2\Vol}\sum_{\vect{k}\vect{p}\vect{p}'ss'}U_2(\vect{k})$\newline$\quad\times c^\dagger_{\vect{p}+\vect{k},s}c^\dagger_{\vect{p}'-\vect{k},s'}c_{\vect{p}'s'}c_{\vect{p}s}$ \\
    $\sum_{i<j}\dfrac{e^2}{4\pi|\vect{r}_i - \vect{r}_j|}$ & $U_2 \to e^2/(4\pi|\vect{r} - \vect{r}'|)$ & $\tfrac1{2\Vol}\sum_{\vect{k}\ne0}v_{\mathrm{C}}(k)\bigl[n_e(\vect{k})^\dagger n_e(\vect{k}) - \op{N}_e\bigr]$ \\
    $-\sum_{iI}\dfrac{Z_{\kappa(I)}e^2}{4\pi|\vect{r}_i - \vect{X}_I|}$ & $\sum_{\gamma\gamma'}\langle\gamma|U_1|\gamma'\rangle c^\dagger_\gamma c_{\gamma'}$, $U_1(\vect{r}) = -\sum_I\dfrac{Z_{\kappa(I)}e^2}{4\pi|\vect{r} - \vect{X}_I|}$ & $-\tfrac1\Vol\sum_{\vect{k}\ne0}v_{\mathrm{C}}(k)\sum_\kappa Z_\kappa\,n_\kappa(\vect{k})\,n_e(\vect{k})^\dagger$ \\
    \bottomrule
  \end{tabular}
\end{table}

\subsection{Change of basis}

A calculation often starts in one basis and ends in another: plane waves for the Coulomb interaction, Bloch states for the band structure (\cref{sec:3.2}), localized Wannier orbitals for interpolation (\cref{sec:3.6}). A new orthonormal basis $|\tilde\gamma\rangle = \sum_\gamma|\gamma\rangle\langle\gamma|\tilde\gamma\rangle$ comes with new operators, defined by the same rule. A creation operator adds a particle in a given orbital, so it expands like the orbital:
\begin{equation}
  c^\dagger_{\tilde\gamma} = \sum_\gamma\langle\gamma|\tilde\gamma\rangle\,c^\dagger_\gamma , \qquad
  c_{\tilde\gamma} = \sum_\gamma\langle\tilde\gamma|\gamma\rangle\,c_\gamma .
  \label{eq:2.3-basis}
\end{equation}
The coefficients form a unitary matrix, so the anticommutators survive,
\begin{equation*}
  \{c_{\tilde\gamma}, c^\dagger_{\tilde\gamma'}\} = \sum_\gamma\langle\tilde\gamma|\gamma\rangle\langle\gamma|\tilde\gamma'\rangle = \delta_{\tilde\gamma\tilde\gamma'} .
\end{equation*}
Inserting \cref{eq:2.3-basis} into \cref{eq:2.3-onebody} and using completeness twice leads to $\sum_{\tilde\gamma\tilde\gamma'}\langle\tilde\gamma|\op{o}|\tilde\gamma'\rangle c^\dagger_{\tilde\gamma}c_{\tilde\gamma'}$: one-body operators look the same in every basis, with the matrix elements of the new one. The field operator of \cref{eq:2.3-field} is the special case in which the new basis is the position eigenbasis, $\langle\gamma|\vect{r}s\rangle = \psi^*_\gamma(\vect{r}, s)$. In the basis that diagonalizes $\op{o}$ the bilinear collapses to $\sum_\gamma o_\gamma c^\dagger_\gamma c_\gamma$. A Slater determinant in the old basis becomes a sum of determinants in the new one (\cref{app:D.2}), which is why the occupation-number language, and not the wave function, is the economical one.

Take two orthonormal orbitals $|A\rangle$ and $|B\rangle$, one on each of two neighboring atoms, and the bonding and antibonding combinations $|\pm\rangle = (|A\rangle \pm |B\rangle)/\sqrt2$, a caricature of the covalent bond of silicon. Two electrons of opposite spin in the bonding orbital form one determinant, $c^\dagger_{+\uparrow}c^\dagger_{+\downarrow}|0\rangle$. By \cref{eq:2.3-basis} it is
\begin{equation*}
  c^\dagger_{+\uparrow}c^\dagger_{+\downarrow}|0\rangle
  = \tfrac12\bigl(c^\dagger_{A\uparrow}c^\dagger_{A\downarrow} + c^\dagger_{B\uparrow}c^\dagger_{B\downarrow}
  + c^\dagger_{A\uparrow}c^\dagger_{B\downarrow} + c^\dagger_{B\uparrow}c^\dagger_{A\downarrow}\bigr)|0\rangle ,
\end{equation*}
four determinants in the atomic basis: two with both electrons on one atom, and two with one electron on each. The same state is simple in one basis and a superposition in the other, and the operators of \cref{eq:2.3-onebody,eq:2.3-twobody} handle both without any change of form.

\subsection{Two-body operators and normal ordering}

The Coulomb repulsion $\tfrac12\sum_{i\ne j}U_2(\vect{r}_i - \vect{r}_j)$, with $U_2(\vect{r}) = e^2/(4\pi r)$, is a sum over pairs, and the condition $i \ne j$ matters: an electron does not repel itself. To see what it becomes, split the pair sum into all pairs minus the diagonal, $\sum_{i\ne j} = \sum_{i,j} - \sum_{i=j}$, and write each sum through the density $\nd_e(\vect{r}) = \sum_i\delta^3(\vect{r} - \vect{r}_i)$:
\begin{equation*}
  \frac12\sum_{i\ne j}U_2(\vect{r}_i - \vect{r}_j)
  = \frac12\int\dd^3r\,\dd^3r'\,U_2(\vect{r} - \vect{r}')\bigl[\nd_e(\vect{r})\nd_e(\vect{r}') - \delta^3(\vect{r} - \vect{r}')\,\nd_e(\vect{r})\bigr] .
\end{equation*}
Each density is now the bilinear $\sum_s\op{\psi}^\dagger_s\op{\psi}_s$. Moving the inner annihilation operator in $\op{\psi}^\dagger_s(\vect{r})\op{\psi}_s(\vect{r})\op{\psi}^\dagger_{s'}(\vect{r}')\op{\psi}_{s'}(\vect{r}')$ to the right with \cref{eq:2.3-field} produces a term $\delta_{ss'}\delta^3(\vect{r} - \vect{r}')\op{\psi}^\dagger_s(\vect{r})\op{\psi}_{s}(\vect{r})$, which cancels the self-interaction term exactly, and leaves, for $\op{O}_2 = \tfrac12\sum_{i\ne j}U_2(\vect{r}_i - \vect{r}_j)$,
\begin{equation}
  \op{O}_2 = \frac12\sum_{ss'}\int\dd^3r\,\dd^3r'\,U_2(\vect{r} - \vect{r}')\,\op{\psi}^\dagger_s(\vect{r})\op{\psi}^\dagger_{s'}(\vect{r}')\op{\psi}_{s'}(\vect{r}')\op{\psi}_s(\vect{r}) .
  \label{eq:2.3-twobody-field}
\end{equation}
All creation operators stand to the left of all annihilation operators. An operator product rearranged into this order, with the sign of the permutation for fermions, is called normal ordered and is written $:\!\cdots\!:$. Thus $\op{O}_2 = \tfrac12\int\!\int U_2\,{:}\nd_e(\vect{r})\nd_e(\vect{r}'){:}$, and normal ordering is what removes the self-interaction. Since a normal-ordered product of two annihilation operators vanishes on any state with fewer than two particles, a single electron feels no Coulomb energy of its own, as it should, and the vacuum has none either. Expanding the field operators in a basis gives the general form of \cref{tab:2-operators},
\begin{equation}
  \op{O}_2 = \frac12\sum_{\gamma_1\gamma_2\gamma_3\gamma_4}\langle\gamma_1\gamma_2|U_2|\gamma_3\gamma_4\rangle\,c^\dagger_{\gamma_1}c^\dagger_{\gamma_2}c_{\gamma_4}c_{\gamma_3} .
  \label{eq:2.3-twobody}
\end{equation}
The order of the last two indices matters: particle 2 is annihilated first and created last, so the two operators of each particle sit symmetrically around the middle. Reversing them flips the sign. \Cref{app:D.2} proves \cref{eq:2.3-twobody} for any two-body operator, spin-dependent or not, by applying the one-body result twice, which is the operator form of the split $\sum_{i,j} - \sum_{i=j}$ used above.

\subsection{Coulomb interaction in momentum space}

The Coulomb potential $e^2/(4\pi r)$ has the Fourier transform (\cref{eq:A.2-vC})
\begin{equation}
  v_{\mathrm{C}}(k) = \frac{e^2}{k^2} = \frac{4\pi\alpha}{k^2} \qquad \text{(Heaviside--Lorentz)} ,
  \label{eq:2.3-vC}
\end{equation}
the $m_V \to 0$ limit of the Yukawa transform of \cref{eq:C.9-Gq}. In Gaussian units the same kernel reads $4\pi e_{\mathrm{G}}^2/k^2$, and in SI units $e_{\mathrm{SI}}^2/(\varepsilon_0^{\mathrm{SI}}k^2)$ (\cref{eq:A.2-map}).

In the box, the inverse transform is a sum, $e^2/(4\pi|\vect{r}|) = \Vol^{-1}\sum_{\vect{k}}v_{\mathrm{C}}(k)\,\ee^{-\ii\vect{k}\cdot\vect{r}}$, and the pair sum factorizes into densities at once. Using the split $\sum_{i\ne j} = \sum_{i,j} - \sum_{i=j}$ again, with $\sum_{i,j}\ee^{-\ii\vect{k}\cdot(\vect{r}_i - \vect{r}_j)} = n_e(-\vect{k})\,n_e(\vect{k})$ and $\sum_{i=j}1 = N_e$,
\begin{equation}
  H_{ee} = \frac12\sum_{i\ne j}\frac{e^2}{4\pi|\vect{r}_i - \vect{r}_j|}
  = \frac{1}{2\Vol}\sum_{\vect{k}}v_{\mathrm{C}}(k)\bigl[n_e(\vect{k})^\dagger n_e(\vect{k}) - \op{N}_e\bigr] .
  \label{eq:2.3-Hee-density}
\end{equation}
The energy of the electrons is a sum over Fourier components of the density fluctuation, each weighted by $v_{\mathrm{C}}(k)$, with the self-energy of each point charge subtracted. Dimensions check: $v_{\mathrm{C}}$ is in $\mathrm{eV}^{-2}$, $1/\Vol$ in $\mathrm{eV}^3$, and the densities are numbers. To get the plane-wave form, insert \cref{eq:2.3-nq} with $n_e(\vect{k})^\dagger = \sum_{\vect{p}'s'}c^\dagger_{\vect{p}'-\vect{k},s'}c_{\vect{p}'s'}$ and normal order. The anticommutator $c_{\vect{p}'s'}c^\dagger_{\vect{p}+\vect{k},s} = \delta_{\vect{p}',\vect{p}+\vect{k}}\delta_{ss'} - c^\dagger_{\vect{p}+\vect{k},s}c_{\vect{p}'s'}$ gives a one-body piece equal to $\op{N}_e$, which cancels the subtraction. In the two-body piece that remains, swapping the two creation operators removes the minus sign:
\begin{equation}
  n_e(\vect{k})^\dagger n_e(\vect{k}) - \op{N}_e = \sum_{\vect{p}\vect{p}'ss'}c^\dagger_{\vect{p}+\vect{k},s}\,c^\dagger_{\vect{p}'-\vect{k},s'}\,c_{\vect{p}'s'}\,c_{\vect{p}s} .
  \label{eq:2.3-normal}
\end{equation}
This is the matrix element of \cref{eq:2.3-twobody} for plane waves: two electrons with momenta $\vect{p}$ and $\vect{p}'$ exchange a momentum $\vect{k}$, each keeps its spin, and the amplitude is $v_{\mathrm{C}}(k)/\Vol$. Without normal ordering, the $\op{N}_e$ term would carry $\Vol^{-1}\sum_{\vect{k}}v_{\mathrm{C}}(k) = \int\dd^3k\,(2\pi)^{-3}e^2/k^2$, which diverges linearly at large $k$: the infinite Coulomb self-energy of a point charge, which has nothing to do with the solid.

The exchange sign shows up as soon as we take a matrix element. Let two electrons start in $c^\dagger_{\vect{p}s}c^\dagger_{\vect{p}'s'}|0\rangle$ and end in $c^\dagger_{\vect{p}+\vect{k},s}c^\dagger_{\vect{p}'-\vect{k},s'}|0\rangle$, with $\vect{k} \ne 0$. The two annihilation operators of \cref{eq:2.3-normal} can remove the initial pair in two ways, with signs $+1$ and $-1$. Each of these pairs with two matchings of the creation operators to the final pair, one at momentum exchange $\vect{k}$ and the other at $\vect{p}' - \vect{p} - \vect{k}$. The four terms add to $2[v_{\mathrm{C}}(k) - \delta_{ss'}v_{\mathrm{C}}(|\vect{p}' - \vect{p} - \vect{k}|)]/\Vol$, and the factor $\tfrac12$ of \cref{eq:2.3-twobody} cancels the 2, which counts the two equivalent labelings of the pair. The result is
\begin{equation}
  \langle\vect{p}+\vect{k}\,s,\ \vect{p}'-\vect{k}\,s'|H_{ee}|\vect{p}s,\ \vect{p}'s'\rangle
  = \frac{1}{\Vol}\Bigl[v_{\mathrm{C}}(k) - \delta_{ss'}\,v_{\mathrm{C}}(|\vect{p}' - \vect{p} - \vect{k}|)\Bigr] .
  \label{eq:2.3-direct-exchange}
\end{equation}
The first term is the direct process: electron $\vect{p}$ takes the momentum $\vect{k}$. The second is the exchange process, in which the electrons end up in the same final orbitals but each in the other's: electron $\vect{p}$ goes to $\vect{p}' - \vect{k}$, so the momentum it takes is $\vect{p}' - \vect{p} - \vect{k}$. Because the two processes lead to the same final state, their amplitudes add, and Fermi statistics attaches a minus sign to the exchange. Electrons with opposite spins, however, can be told apart by their spins, and for them the exchange term is absent. The exchange term lowers the energy of parallel spins, and it is the origin of the exchange energy of \cref{sec:3.4} and of the exchange coupling of magnetic order (\cref{sec:7.1}).

The same counting gives the Coulomb energy of a single determinant, the starting point of the Hartree--Fock method. In $\langle\Phi_0|c^\dagger_{\gamma_1}c^\dagger_{\gamma_2}c_{\gamma_4}c_{\gamma_3}|\Phi_0\rangle$ the annihilation operators must remove two occupied orbitals and the creation operators must put the same two back, either in the same places or swapped. For orthonormal orbitals \cref{eq:2.3-twobody} then reduces to
\begin{equation*}
  \langle\Phi_0|\op{O}_2|\Phi_0\rangle = \frac12\sum_{\gamma,\gamma'\ \mathrm{occupied}}\Bigl[\langle\gamma\gamma'|U_2|\gamma\gamma'\rangle - \langle\gamma\gamma'|U_2|\gamma'\gamma\rangle\Bigr],
\end{equation*}
a direct (Hartree) term, the classical electrostatic energy of the charge distribution, minus an exchange (Fock) term. The diagonal $\gamma = \gamma'$ cancels between the two: no orbital interacts with itself. \Cref{sec:3.4} minimizes this energy and then improves on it.

\subsection{Coulomb Hamiltonian of a solid}

\Cref{eq:2.3-Hee-density} has one more infinity, at $\vect{k} = 0$, where $v_{\mathrm{C}}(0)$ diverges. It is the energy of a uniform charge spread over the whole box, which grows faster than the volume, and in a neutral system it is cancelled by the positive charge. The simplest model of a metal or of a valence electron gas, jellium, replaces the nuclei by a uniform positive background. Jellium (used again in \cref{sec:2.5,sec:4.5}) shows how the cancellation works, in \cref{we:2-jellium} below.

How large is the Coulomb term? A rough answer comes from the mean spacing $d = \nd_e^{-1/3}$ of the electrons. For the four valence electrons per atom of silicon, $\nd_e = \SI{2.00e23}{cm^{-3}}$ (\cref{app:A.2}), and $d = \SI{1.71}{\angstrom} = 3.2\,a_0$, with $a_0 = 1/(\alpha m_e)$ the Bohr radius (\cref{sec:2.2}). The Coulomb energy of two neighbors, $\alpha/d$, is then about \SI{8.4}{eV}. The kinetic energy must respect the Pauli principle, and for a free gas at this density its mean is $\tfrac35E_{\mathrm{F}} = \SI{7.5}{eV}$ (\cref{sec:2.5}), about the same. The field measures the ratio by $r_s = (3/4\pi\nd_e)^{1/3}/a_0$, the radius per electron in units of $a_0$: the Coulomb energy scales as $1/r_s$ and the kinetic energy as $1/r_s^2$, and their ratio grows as $r_s$. For silicon's valence density $r_s = 2.0$. The interaction is therefore not a small perturbation of free electrons. It is tamed instead by screening (\cref{sec:4.4}) and by the self-consistent fields of \cref{sec:3.4}, and all of them start from the Coulomb Hamiltonian of \cref{we:2-jellium}.

\begin{workedexample}{The Coulomb Hamiltonian of jellium and of a solid}\label{we:2-jellium}
Take $N_e$ electrons in the box with a uniform positive background of number density $N_e/\Vol$ and charge $+e$ per unit, so the system is neutral. To give the $\vect{k} = 0$ terms a meaning, screen the Coulomb potential at a large length $1/k_{\mathrm{reg}}$, so $e^2/(4\pi r) \to e^2\ee^{-k_{\mathrm{reg}}r}/(4\pi r)$ and $v_{\mathrm{C}}(k) \to e^2/(k^2 + k_{\mathrm{reg}}^2)$ by \cref{eq:C.9-Gq}. We take $\Vol \to \infty$ at fixed $N_e/\Vol$ first, with $\Vol^{1/3} \gg 1/k_{\mathrm{reg}}$, and $k_{\mathrm{reg}} \to 0$ at the end. The integral of the screened potential over all space is $v_{\mathrm{C}}(0) = e^2/k_{\mathrm{reg}}^2$. The uniform background has a Fourier component only at $\vect{k} = 0$, where three terms meet:
\begin{itemize}
  \item the background with itself, $+\tfrac12(N_e/\Vol)^2\,\Vol\,v_{\mathrm{C}}(0) = +N_e^2\,v_{\mathrm{C}}(0)/2\Vol$;
  \item the electrons with the background, $-(N_e/\Vol)\sum_i\int\dd^3r\,e^2\ee^{-k_{\mathrm{reg}}|\vect{r}_i - \vect{r}|}/(4\pi|\vect{r}_i - \vect{r}|) = -N_e^2\,v_{\mathrm{C}}(0)/\Vol$;
  \item the $\vect{k} = 0$ term of \cref{eq:2.3-Hee-density}, $+v_{\mathrm{C}}(0)\,(N_e^2 - N_e)/2\Vol$, since $n_e(\vect{0}) = \op{N}_e = N_e$.
\end{itemize}
The terms in $N_e^2$, each of order $N_e\cdot(N_e/\Vol)/k_{\mathrm{reg}}^2$ and so proportional to the volume, cancel exactly. What is left is $-N_e\,v_{\mathrm{C}}(0)/2\Vol = -N_e e^2/(2Vk_{\mathrm{reg}}^2)$, and per electron it vanishes as $\Vol \to \infty$. Then $k_{\mathrm{reg}} \to 0$ restores $v_{\mathrm{C}}(k) = e^2/k^2$ at every $\vect{k} \ne 0$. With the kinetic energy of \cref{tab:2-operators} and \cref{eq:2.3-normal}, the jellium Hamiltonian is~\cite[Chap.~2]{bruus2004many}
\begin{equation}
  H_{\mathrm{jel}} = \sum_{\vect{k}s}\frac{k^2}{2m_e}\,c^\dagger_{\vect{k}s}c_{\vect{k}s}
  + \frac{1}{2\Vol}\sum_{\vect{k}\ne0}\sum_{\vect{p}\vect{p}'ss'}v_{\mathrm{C}}(k)\,c^\dagger_{\vect{p}+\vect{k},s}\,c^\dagger_{\vect{p}'-\vect{k},s'}\,c_{\vect{p}'s'}\,c_{\vect{p}s} .
  \label{eq:2.3-jellium}
\end{equation}
The background has disappeared, apart from the instruction to drop $\vect{k} = 0$. The order of limits is part of the answer: with $k_{\mathrm{reg}} \to 0$ first, the screening length would exceed the box, and the three terms would each be infinite.

A solid replaces the background by point nuclei. Nucleus $I$ sits at $\vect{X}_I$, belongs to a species $\kappa(I)$ with charge number $Z_\kappa$ and mass $M_\kappa$ (for silicon one species, $Z_\kappa = 14$; species are treated systematically in \cref{sec:5.2}), and the species density is $n_\kappa(\vect{k}) = \sum_{I\in\kappa}\ee^{\ii\vect{k}\cdot\vect{X}_I}$. Every pair of point charges interacts through the product of their charge numbers, $-1$ for an electron and $Z_\kappa$ for a nucleus, times $e^2/(4\pi r)$ at separation $r$. Splitting the pair sum as in \cref{eq:2.3-Hee-density} leads to
\begin{equation*}
  H_{\mathrm{Coul}} = \frac{1}{2\Vol}\sum_{\vect{k}}v_{\mathrm{C}}(k)\Bigl[n_Z(\vect{k})^\dagger n_Z(\vect{k}) - N_e - \sum_\kappa Z_\kappa^2N_\kappa\Bigr], \qquad
  n_Z(\vect{k}) = -n_e(\vect{k}) + \sum_\kappa Z_\kappa\,n_\kappa(\vect{k}),
\end{equation*}
where $N_\kappa$ counts the nuclei of species $\kappa$. Neutrality, $\sum_\kappa Z_\kappa N_\kappa = N_e$, now holds as an operator identity, $n_Z(\vect{0}) = 0$. The $\vect{k} = 0$ term reduces to the self-energy subtraction, $-v_{\mathrm{C}}(0)(N_e + \sum_\kappa Z_\kappa^2N_\kappa)/2\Vol$, which vanishes per particle in the same limit as before. The nuclei are their own neutralizing background. Expanding $n_Z^\dagger n_Z$ leads to three pieces. The electron--electron piece is \cref{eq:2.3-normal}. The cross terms $-Z_\kappa[n_e(\vect{k})^\dagger n_\kappa(\vect{k}) + n_\kappa(\vect{k})^\dagger n_e(\vect{k})]$ are equal after $\vect{k} \to -\vect{k}$, since $v_{\mathrm{C}}$ is even and electron and nuclear coordinates commute, and add to twice the first. The nuclear piece keeps its own self-energy subtraction.
\end{workedexample}

Collecting the pieces, with the nuclei left in first quantization because their quantization belongs to the lattice (\cref{sec:5.3}):
\begin{keyresult}[Coulomb Hamiltonian of electrons and point nuclei]
\begin{multline}
  H_C = \sum_{\vect{k}s}\frac{k^2}{2m_e}c^\dagger_{\vect{k}s}c_{\vect{k}s} - \sum_I\frac{\nabla_I^2}{2M_{\kappa(I)}}
  + \frac{1}{2\Vol}\sum_{\vect{k}\ne0}v_{\mathrm{C}}(k)\Bigl[\sum_{\vect{p}\vect{p}'ss'}c^\dagger_{\vect{p}+\vect{k},s}c^\dagger_{\vect{p}'-\vect{k},s'}c_{\vect{p}'s'}c_{\vect{p}s}\\
  - 2\sum_\kappa Z_\kappa\,n_\kappa(\vect{k})\sum_{\vect{p}s}c^\dagger_{\vect{p}-\vect{k},s}c_{\vect{p}s}
  + \sum_{\kappa\kappa'}Z_\kappa Z_{\kappa'}\bigl(n_\kappa(\vect{k})^\dagger n_{\kappa'}(\vect{k}) - \delta_{\kappa\kappa'}N_\kappa\bigr)\Bigr] .
  \label{eq:2.3-HC}
\end{multline}
The three terms in the bracket are the electron--electron, electron--nucleus and nucleus--nucleus interactions; $\nabla_I$ acts on $\vect{X}_I$, and $\vect{k} = 0$ is removed by neutrality.
\end{keyresult}
The powers of $\Vol$ in \cref{eq:2.3-HC} are a useful check, in the spirit of \cref{sec:2.1}. The energy of a large system must be extensive, proportional to $\Vol$ at fixed density. In the electron--electron term, take the expectation value in a determinant of plane waves. Only terms that put both electrons back where they started contribute. The direct one has $\vect{k} = 0$ and is gone with the background; the exchange one, the diagonal case of \cref{eq:2.3-direct-exchange} with $\vect{k} = \vect{p}' - \vect{p}$, leaves $-(1/2\Vol)\sum_{\vect{p}\ne\vect{p}'}v_{\mathrm{C}}(|\vect{p} - \vect{p}'|)$ for each spin, summed over occupied momenta. By \cref{eq:2.1-box} each sum over momenta brings a factor $\Vol$, and we see that the total is $\Vol^2/\Vol = \Vol$ times an integral over the occupied region: extensive, as it must be. The same count applied to the subtracted self-energy, $(1/2\Vol)\sum_{\vect{k}}v_{\mathrm{C}}(k)N_e$, yields $N_e$ times a divergent integral, the reason it is removed rather than kept. A stray power of $\Vol$ in a Coulomb expression nearly always means a misplaced $\vect{k} = 0$ term or a forgotten box normalization.

The electron--nucleus term is the one-body potential $\sum_iU_1(\vect{r}_i)$ of \cref{tab:2-operators} with $U_1(\vect{k}) = -v_{\mathrm{C}}(k)\sum_\kappa Z_\kappa n_\kappa(\vect{k})$. As a check, its matrix element for an electron going from $\vect{p}$ to $\vect{p} - \vect{k}$,
\begin{equation*}
  \langle\vect{p} - \vect{k}|U_1|\vect{p}\rangle = -\frac{1}{\Vol}\sum_IZ_{\kappa(I)}\,v_{\mathrm{C}}(k)\,\ee^{\ii\vect{k}\cdot\vect{X}_I} ,
\end{equation*}
is what \cref{eq:A.3-vertex} gives for the potential $-\sum_IZ_{\kappa(I)}e^2/(4\pi|\vect{r} - \vect{X}_I|)$, with the nuclei in the role of the dark matter. When the nuclei are held at fixed positions, $n_\kappa(\vect{k})$ is a number, and for a perfect crystal it is nonzero only at reciprocal-lattice vectors, which leads to the periodic potential of the band structure (\cref{sec:3.1,sec:3.2}).

Momentum conservation reads off \cref{eq:2.3-HC} directly. In every electron--electron term the electrons gain $\vect{k}$ and lose $\vect{k}$, and the total electron momentum is conserved by the Coulomb repulsion. The electron--nucleus term moves electron momentum by $-\vect{k}$ and multiplies by $n_\kappa(\vect{k})$, which, for mobile nuclei, shifts the nuclear momenta by $+\vect{k}$: the total momentum of electrons and nuclei is conserved. With the nuclei clamped in a crystal, the electrons can gain or lose only reciprocal-lattice vectors, and their momentum is conserved only modulo these. This is the origin of crystal momentum (\cref{sec:3.1}).

\Cref{eq:2.3-HC} is the starting point of everything the book computes about the target. It is nonrelativistic and treats the nuclei as point charges; spin--orbit coupling and the relativistic corrections of core electrons are added where they matter (\cref{sec:3.3,sec:3.5}). \Cref{sec:5.1} separates the motion of the nuclei from that of the electrons, and the dark-matter coupling enters only through $n_e(\vect{q})$ for electrons, or $n_\kappa(\vect{q})$ for nuclei. No one, however, solves \cref{eq:2.3-HC} for $10^{26}$ electrons as it stands. Density-functional theory replaces the electron--electron term by an effective one-body potential (\cref{sec:3.4}), and the codes that do this in practice are described in \cref{app:E.1}.
